\section{多维傅里叶变换}\label{sec:Multi_Fourier}

这一节中，我们将傅里叶变换推广到高维。由于时间总是一维的，高维傅里叶变换往往指空间域，因此我们将用$\mathbf{x}=(x_1,x_2,\cdots,x_n)$表示n维的空间变量，$\boldsymbol{\xi}=(\xi_1,\xi_2,\cdots,\xi_n)$表示n维的空间频率，$\boldsymbol{\omega}=(\omega_1,\omega_2,\cdots,\omega_n)$表示n维的空间角频率。多维傅里叶变换在图像处理、计算机视觉等领域有着重要的应用。

\subsection{多维傅里叶变换的定义}\label{sub:multi_fourier}

\begin{definition}[n维傅里叶变换]
设$f\in L^1(\mathbb{R}^n)$，则其n维傅里叶变换定义为：
\begin{lralign}
    \mathcal{F} f(\boldsymbol{\xi})&=\int_{\mathbb{R}^n}f(\mathbf{x})\mathrm{e}^{-2\pi \mathrm{i}\mathbf{x}\cdot\boldsymbol{\xi}}\,d\mathbf{x}
\switch
    \mathcal{F} f(\boldsymbol{\omega})&=\int_{\mathbb{R}^n}f(\mathbf{x})\mathrm{e}^{-\mathrm{i}\mathbf{x}\cdot\boldsymbol{\omega}}\,d\mathbf{x}
\end{lralign}
我们仍记$f\overset{\mathcal{F} }{\longleftrightarrow}F$并称$f$与$F$构成\textbf{傅里叶变换对}。
\end{definition}

类似一维傅里叶变换，我们来验证$f\in L^1(\mathbb{R}^n)$时，其多维傅里叶变换是有界的、一致连续的。
\begin{proposition}
    设$f\in L^1(\mathbb{R}^n)$，则其n维傅里叶变换$\mathcal{F} f$是有界的、一致连续的。
\end{proposition}
\begin{proof}
    有界性：由定义，有
\begin{lralign}
    |\mathcal{F} f(\boldsymbol{\xi})|&=\left|\int_{\mathbb{R}^n}f(\mathbf{x})\mathrm{e}^{-2\pi \mathrm{i}\mathbf{x}\cdot\boldsymbol{\xi}}\,d\mathbf{x}\right|\\
    &\leq \int_{\mathbb{R}^n}|f(\mathbf{x})||\mathrm{e}^{-2\pi \mathrm{i}\mathbf{x}\cdot\boldsymbol{\xi}}|\,d\mathbf{x}\\
    &= \int_{\mathbb{R}^n}|f(\mathbf{x})|\,d\mathbf{x}=\|f\|_1<\infty
\switch
    |\mathcal{F} f(\boldsymbol{\omega})|&=\left|\int_{\mathbb{R}^n}f(\mathbf{x})\mathrm{e}^{-\mathrm{i}\mathbf{x}\cdot\boldsymbol{\omega}}\,d\mathbf{x}\right|\\
    &\leq \int_{\mathbb{R}^n}|f(\mathbf{x})||\mathrm{e}^{-\mathrm{i}\mathbf{x}\cdot\boldsymbol{\omega}}|\,d\mathbf{x}\\
    &= \int_{\mathbb{R}^n}|f(\mathbf{x})|\,d\mathbf{x}=\|f\|_1<\infty
\end{lralign}
一致连续性：设$\boldsymbol{\xi},\boldsymbol{\eta}\in\mathbb{R}^n$，则
\begin{lralign}
    &|\mathcal{F} f(\boldsymbol{\xi})-\mathcal{F} f(\boldsymbol{\eta})|\\
    =&\left|\int_{\mathbb{R}^n}f(\mathbf{x})\left(\mathrm{e}^{-2\pi \mathrm{i}\mathbf{x}\cdot\boldsymbol{\xi}}-\mathrm{e}^{-2\pi \mathrm{i}\mathbf{x}\cdot\boldsymbol{\eta}}\right)\,d\mathbf{x}\right|\\
    \leq& \int_{\mathbb{R}^n}|f(\mathbf{x})|\left|\mathrm{e}^{-2\pi \mathrm{i}\mathbf{x}\cdot\boldsymbol{\xi}}-\mathrm{e}^{-2\pi \mathrm{i}\mathbf{x}\cdot\boldsymbol{\eta}}\right|\,d\mathbf{x}
\switch
    &|\mathcal{F} f(\boldsymbol{\omega})-\mathcal{F} f(\boldsymbol{\nu})|\\
    =&\left|\int_{\mathbb{R}^n}f(\mathbf{x})\left(\mathrm{e}^{-\mathrm{i}\mathbf{x}\cdot\boldsymbol{\omega}}-\mathrm{e}^{-\mathrm{i}\mathbf{x}\cdot\boldsymbol{\nu}}\right)\,d\mathbf{x}\right|\\
    \leq& \int_{\mathbb{R}^n}|f(\mathbf{x})|\left|\mathrm{e}^{-\mathrm{i}\mathbf{x}\cdot\boldsymbol{\omega}}-\mathrm{e}^{-\mathrm{i}\mathbf{x}\cdot\boldsymbol{\nu}}\right|\,d\mathbf{x}
\end{lralign}
由于$\mathrm{e}^{-2\pi \mathrm{i}\mathbf{x}\cdot\boldsymbol{\xi}}$与$\mathrm{e}^{-\mathrm{i}\mathbf{x}\cdot\boldsymbol{\omega}}$均为一致连续函数，$\forall\epsilon>0,\exists\delta>0$，使得\begin{lralign}
    &|\boldsymbol{\xi}-\boldsymbol{\eta}|<\delta\\
    \implies &|\mathrm{e}^{-2\pi \mathrm{i}\mathbf{x}\cdot\boldsymbol{\xi}}-\mathrm{e}^{-2\pi \mathrm{i}\mathbf{x}\cdot\boldsymbol{\eta}}|<\frac{\epsilon}{\|f\|_1},\forall \mathbf{x}\in\mathbb{R}^n\\
    \implies &|\mathcal{F} f(\boldsymbol{\xi})-\mathcal{F} f(\boldsymbol{\eta})|<\epsilon
\switch
    &|\boldsymbol{\omega}-\boldsymbol{\nu}|<\delta\\
    \implies &|\mathrm{e}^{-\mathrm{i}\mathbf{x}\cdot\boldsymbol{\omega}}-\mathrm{e}^{-\mathrm{i}\mathbf{x}\cdot\boldsymbol{\nu}}|<\frac{\epsilon}{\|f\|_1},\forall \mathbf{x}\in\mathbb{R}^n\\
    \implies &|\mathcal{F} f(\boldsymbol{\omega})-\mathcal{F} f(\boldsymbol{\nu})|<\epsilon
\end{lralign}
\end{proof}

此外，黎曼-勒贝格引理仍然成立：
\begin{claim}[n维黎曼-勒贝格引理]
    设$f\in L^1(\mathbb{R}^n)$，则其n维傅里叶变换$\mathcal{F} f$满足
    \begin{lralign}
    \lim_{|\boldsymbol{\xi}|\to\infty}\mathcal{F} f(\boldsymbol{\xi})&=0
\switch
    \lim_{|\boldsymbol{\omega}|\to\infty}\mathcal{F} f(\boldsymbol{\omega})&=0
\end{lralign}
\end{claim}

不难想到，傅里叶逆变换为：
\begin{definition}[n维傅里叶逆变换]
    设$F\in L^1(\mathbb{R}^n)$，则其n维傅里叶逆变换为：
    \begin{lralign}
    \mathcal{F} ^{-1} F(\mathbf{x})&=\int_{\mathbb{R}^n}F(\boldsymbol{\xi})\mathrm{e}^{2\pi \mathrm{i}\mathbf{x}\cdot\boldsymbol{\xi}}\,d\boldsymbol{\xi}
\switch
    \mathcal{F} ^{-1} F(\mathbf{x})&=\frac{1}{(2\pi)^n}\int_{\mathbb{R}^n}F(\boldsymbol{\omega})\mathrm{e}^{\mathrm{i}\mathbf{x}\cdot\boldsymbol{\omega}}\,d\boldsymbol{\omega}
\end{lralign}
\end{definition}

傅里叶反演公式仍然成立，因此n维傅里叶变换也是一一的，从而前文中傅里叶变换对的记号$f\overset{\mathcal{F} }{\longleftrightarrow}F$是合理的。
\begin{claim}[n维傅里叶反演公式]
    设$f\in L^1(\mathbb{R}^n)$且$\mathcal{F} f\in L^1(\mathbb{R}^n)$，则
    \begin{lralign}
    f(\mathbf{x})&=\int_{\mathbb{R}^n}\mathcal{F} f(\boldsymbol{\xi})\mathrm{e}^{2\pi \mathrm{i}\mathbf{x}\cdot\boldsymbol{\xi}}\,d\boldsymbol{\xi}
\switch
    f(\mathbf{x})&=\frac{1}{(2\pi)^n}\int_{\mathbb{R}^n}\mathcal{F} f(\boldsymbol{\omega})\mathrm{e}^{\mathrm{i}\mathbf{x}\cdot\boldsymbol{\omega}}\,d\boldsymbol{\omega}
\end{lralign}
\end{claim}
\begin{remark}
    多维情况下很难讲清楚什么是“左右极限平均值”，一种处理方式是通过球平均在半径趋于0时的极限来替代$\dfrac{f(t+)+f(t-)}{2}$，但这样比较复杂，因此我们不妨说$\mathcal{F} f$的傅里叶逆变换几乎处处收敛到$f$，或者要求$f\in C(\mathbb{R}^n)$。
\end{remark}

为了说明空间角频率形式的傅里叶逆变换的系数$\dfrac{1}{(2\pi)^n}$的由来，我们可以令$\boldsymbol{\omega}=2\pi\boldsymbol{\xi}$，则$d\boldsymbol{\omega}=(2\pi)^n d\boldsymbol{\xi}$，以上两种傅里叶逆变换的公式是等价的。

下面给出几个例子来建立多维傅里叶变换的直观。
\begin{definition}[可分函数]
    $f(\mathbf{x}):\mathbb{R}^n\to\mathbb{C}$是一个可分函数，即如果存在$n$个一维函数$f_i(x_i):\mathbb{R}\to\mathbb{C},\mathrm{i}=1,2,\cdots,n$，使得
    \begin{align*}
        f(\mathbf{x})=f_1(x_1)f_2(x_2)\cdots f_n(x_n)
    \end{align*}
\end{definition}
$f(\mathbf{x})$的傅里叶变换为（以频率为例）：

{\color{blue}
\begin{align*}
\mathcal{F} f(\boldsymbol{\xi})&=\int_{\mathbb{R}^n}f(\mathbf{x})\mathrm{e}^{-2\pi \mathrm{i} \mathbf{x}\cdot\boldsymbol{\xi}}\,d\mathbf{x}\\
&=\int_{-\infty}^{\infty}\cdots \int_{-\infty}^{\infty}f_1(x_1)f_2(x_2)\cdots f_n(x_n)\mathrm{e}^{-2\pi \mathrm{i} (x_1\xi_1+x_2\xi_2+\cdots +x_n\xi_n)}\,dx_1 dx_2\cdots dx_n\\
&=\left(\int_{-\infty}^{\infty}f_1(x_1)\mathrm{e}^{-2\pi \mathrm{i} x_1\xi_1}\,dx_1\right)\left(\int_{-\infty}^{\infty}f_2(x_2)\mathrm{e}^{-2\pi \mathrm{i} x_2\xi_2}\,dx_2\right)\cdots \left(\int_{-\infty}^{\infty}f_n(x_n)\mathrm{e}^{-2\pi \mathrm{i} x_n\xi_n}\,dx_n\right)\\
&=\mathcal{F} f_1(\xi_1)\mathcal{F} f_2(\xi_2)\cdots \mathcal{F} f_n(\xi_n)
\end{align*}}
%{\color{red}
%\begin{align*}
%\mathcal{F} f(\boldsymbol{\omega})&=\int_{\mathbb{R}^n}f(\mathbf{x})\mathrm{e}^{-\mathrm{i} \mathbf{x}\cdot\boldsymbol{\omega}}\,d\mathbf{x}\\
%&=\int_{-\infty}^{\infty}\cdots \int_{-\infty}^{\infty}f_1(x_1)f_2(x_2)\cdots f_n(x_n)\mathrm{e}^{-\mathrm{i} (x_1\omega_1+x_2\omega_2+\cdots +x_n\omega_n)}\,dx_1 dx_2\cdots dx_n\\  
%&=\left(\int_{-\infty}^{\infty}f_1(x_1)\mathrm{e}^{-\mathrm{i} x_1\omega_1}\,dx_1\right)\left(\int_{-\infty}^{\infty}f_2(x_2)\mathrm{e}^{-\mathrm{i} x_2\omega_2}\,dx_2\right)\cdots \left(\int_{-\infty}^{\infty}f_n(x_n)\mathrm{e}^{-\mathrm{i} x_n\omega_n}\,dx_n\right)\\
%&=\mathcal{F} f_1(\omega_1)\mathcal{F} f_2(\omega_2)\cdots \mathcal{F} f_n(\omega_n)
%\end{align*}}
也就是说，空间域上可分的函数，其傅里叶变换也是可分的，且每个分量的傅里叶变换正是对应一维函数的傅里叶变换。

\begin{example}[高斯函数]
    高斯函数是可分的：
    \begin{align*}
        f(\mathbf{x})=\mathrm{e}^{-\pi(x_1^2+x_2^2+\cdots +x_n^2)}=\mathrm{e}^{-\pi x_1^2}\mathrm{e}^{-\pi x_2^2}\cdots \mathrm{e}^{-\pi x_n^2}
    \end{align*}
    我们仍考虑这个特殊的高斯函数，并指出其他高斯函数可以通过下一小节中介绍的运算性质得到。一维高斯函数的傅里叶变换为：
    \begin{lralign}
    \mathcal{F} \left(\mathrm{e}^{-\pi x^2}\right)(\xi)&=\mathrm{e}^{-\pi \xi^2}
\switch
    \mathcal{F} \left(\mathrm{e}^{-\pi x^2}\right)(\omega)&=\mathrm{e}^{-\frac{\omega^2}{4\pi}}
    \end{lralign}
    因此，n维高斯函数的傅里叶变换仍为高斯函数：
    \begin{lralign}
    \mathcal{F} f(\boldsymbol{\xi})=\mathrm{e}^{-\pi(\xi_1^2+\xi_2^2+\cdots +\xi_n^2)}=\mathrm{e}^{-\pi |\boldsymbol{\xi}|^2}\quad
\switch
    \mathcal{F} f(\boldsymbol{\omega})=\mathrm{e}^{-(\omega_1^2+\omega_2^2+\cdots +\omega_n^2)/4\pi}=\mathrm{e}^{-|\boldsymbol{\omega}|^2/4\pi}\quad
    \end{lralign}
    \end{example}
    \begin{figure}[H]
        \centering
        \includegraphics[width=0.5\textwidth]{n_gauss.jpeg}
        \caption{二维高斯函数的图像}
    \end{figure}

\begin{example}[矩形函数]
    n维空间中的矩形函数是n维“长方体”的特征函数。由于这种长方体是区间(Interval)的笛卡尔积，是n维空间中区间的推广，我们仍称之为区间并用$I$表示：
    \begin{align*}
        I=[a_1,b_1]\times[a_2,b_2]\times\cdots \times[a_n,b_n]
    \end{align*}
    n维矩形函数记为$\chi_I$。简洁起见，这里计算一个特殊的矩形函数$\chi_I$，其中$I=[-T/2,T/2]^n$。已知一维矩形函数的傅里叶变换为：
    \begin{lralign}
    \mathcal{F} \left(\chi_{[-T/2,T/2]}\right)(\xi)&=T\text{sinc}(T\xi)
\switch
    \mathcal{F} \left(\chi_{[-T/2,T/2]}\right)(\omega)&=T\text{Sa}\left(\frac{T\omega}{2}\right)
    \end{lralign}
    因此，n维矩形函数的傅里叶变换为：
    \begin{lralign}
    \mathcal{F} \chi_I(\boldsymbol{\xi})&=T^n \prod_{j=1}^n \text{sinc}(T\xi_j)
\switch
    \mathcal{F} \chi_I(\boldsymbol{\omega})&=T^n \prod_{j=1}^n \text{Sa}\left(\frac{T\omega_j}{2}\right)
    \end{lralign}
\end{example}

\begin{wrapfigure}{r}{0.5\textwidth}
    \centering
    \vspace{-20pt}
    \includegraphics[width=0.48\textwidth]{zero_phase.jpeg}
    \caption{二维空间中$\mathrm{e}^{2\pi \mathrm{i} \mathbf{x}\cdot\boldsymbol{\xi}}$的等相面示意图}
    \vspace{4em}
\end{wrapfigure}
下面考察$\mathrm{e}^{2\pi \mathrm{i} \mathbf{x}\cdot\boldsymbol{\xi}}$和$\mathrm{e}^{\mathrm{i} \mathbf{x}\cdot\boldsymbol{\omega}}$的性质。以$\mathrm{e}^{2\pi \mathrm{i} \mathbf{x}\cdot\boldsymbol{\xi}}$为例，对于给定的$\mathbf{\xi}$，令$$\boldsymbol{x\cdot\xi}=k(k\in\mathbb{Z})$$
则$$\mathrm{e}^{2\pi \mathrm{i} \mathbf{x}\cdot\boldsymbol{\xi}}=1$$
这时称$\mathrm{e}^{2\pi \mathrm{i} \mathbf{x}\cdot\boldsymbol{\xi}}$是\textbf{零相}(zero phase)的。
这些点构成了过原点且法向量为$\boldsymbol{\xi_0}$的超平面族，相邻零相面之间的距离为$1/|\boldsymbol{\xi}|$。类似地，可以定义\textbf{等相面}(constant phase surface)为
$$\left\{\mathbf{x}\in\mathbb{R}^n \mid \mathbf{x}\cdot\boldsymbol{\xi}=k+\theta,k\in\mathbb{Z},\theta\in[0,1)\right\}$$
的超平面族，法向量为$\boldsymbol{\xi_0}$，相邻等相面距离为$1/|\boldsymbol{\xi}|$。

\subsection{多维傅里叶变换的运算性质}\label{sub:operation_of_multi_fourier}

多维傅里叶变换的运算性质相较一维比较丰富，本小节中将分条列举这些性质。由于

\begin{proposition}[对偶性]
    设$f(\mathbf{x}):\mathbb{R}^n\to\mathbb{C}$，则有：
    \begin{lralign}
    &\mathcal{F} f^- =\mathcal{F} (f^-)=\mathcal{F} ^{-1} f\\
    &\mathcal{F} \left(\mathcal{F} f\right)=f
\switch
    &\mathcal{F} f^-=(2\pi)^n \mathcal{F} ^{-1} f\\
    &\mathcal{F} \left(\mathcal{F} f\right)=(2\pi)^n f^-
    \end{lralign}
    进一步，如果$f$是实函数，则有：
    \begin{lralign}
    \mathcal{F} f^-&=\overline{\mathcal{F} f}
\switch
    \mathcal{F} f^-&=(2\pi)^n \overline{\mathcal{F} f}
    \end{lralign}
\end{proposition}
\begin{proof}
    类似一维的情况（见\ref{sub:fourier_operation_properties}），带入定义可立即得到：
    \begin{lralign}
        \mathcal{F} f^-=\mathcal{F} (f^-)=\mathcal{F} ^{-1} f
    \switch
        \mathcal{F} f^-=(2\pi)^n \mathcal{F} ^{-1} f
    \end{lralign}
    同时取傅里叶变换，即得：
    \begin{lralign}
        \mathcal{F} \left(\mathcal{F} f\right)=f^-
\switch
        \mathcal{F} \left(\mathcal{F} f\right)=(2\pi)^n f^-
    \end{lralign}
    如果$f$是实函数，则$f=\overline{f}$，因此：
    \begin{lralign}
    \mathcal{F} f^-(\mathbf{\xi})&=\int_{\mathbb{R}^n}f(\mathbf{x})\mathrm{e}^{2\pi \mathrm{i} \mathbf{x}\cdot\boldsymbol{\xi}}\,d\mathbf{x}\\
    &=\overline{\int_{\mathbb{R}^n}f(\mathbf{x})\mathrm{e}^{-2\pi \mathrm{i} \mathbf{x}\cdot\boldsymbol{\xi}}\,d\mathbf{x}}\\
    &=\overline{\mathcal{F} f(\boldsymbol{\xi})}
\switch
    \mathcal{F} f^-(\mathbf{\omega})&=\int_{\mathbb{R}^n}f(-\mathbf{x})\mathrm{e}^{-\mathrm{i} \mathbf{x}\cdot\mathbf{\omega}}\,d\mathbf{x}\\
    &=\overline{\int_{\mathbb{R}^n}f(\mathbf{x})\mathrm{e}^{\mathrm{i} \mathbf{x}\cdot\mathbf{\omega}}\,d\mathbf{x}}\\
    &=(2\pi)^n \overline{\mathcal{F} f(\mathbf{\omega})}
    \end{lralign}
\end{proof}

\begin{corollary}[对称性]
    可仿照一维函数定义函数的奇偶性：如果$f(\mathbf{x})=f(-\mathbf{x})$，则称$f$为偶函数；如果$f(\mathbf{x})=-f(-\mathbf{x})$，则称$f$为奇函数。$\mathcal{F} f$与$f$具有相同的奇偶性；$f$是实函数时，如果$f$为偶函数，则$\mathcal{F} f$为实函数；如果$f$为奇函数，则$\mathcal{F} f$为纯虚函数。
\end{corollary}

\begin{proposition}[线性性]
    设$f(\mathbf{x}),g(\mathbf{x}):\mathbb{R}^n\to\mathbb{C}$，则有：
    \begin{lralign}
    \mathcal{F} [af(\mathbf{x})+bg(\mathbf{x})](\boldsymbol{\xi})&=a\mathcal{F} f(\boldsymbol{\xi})+b\mathcal{F} g(\boldsymbol{\xi})    
\switch
    \mathcal{F} [af(\mathbf{x})+bg(\mathbf{x})](\boldsymbol{\omega})&=a\mathcal{F} f(\boldsymbol{\omega})+b\mathcal{F} g(\boldsymbol{\omega})
    \end{lralign}
\end{proposition}
这个性质是显然的，它直接得自积分的线性性。

\begin{proposition}[平移定理]
    设$f(\mathbf{x}):\mathbb{R}^n\to\mathbb{C}$，时域上：
    \begin{lralign}
    \mathcal{F} \left[f(\mathbf{x}-\mathbf{a})\right](\boldsymbol{\xi})&=\mathrm{e}^{-2\pi \mathrm{i} \mathbf{a}\cdot\boldsymbol{\xi}}\mathcal{F} f(\boldsymbol{\xi})
\switch
    \mathcal{F} \left[f(\mathbf{x}-\mathbf{a})\right](\boldsymbol{\omega})&=\mathrm{e}^{-\mathrm{i} \mathbf{a}\cdot\boldsymbol{\omega}}\mathcal{F} f(\boldsymbol{\omega})
    \end{lralign}
    频域上：
    \begin{lralign}
    \mathcal{F} \left[\mathrm{e}^{2\pi \mathrm{i} \mathbf{a}\cdot\boldsymbol{\xi}}f(\mathbf{x})\right](\boldsymbol{\xi})&=f(\boldsymbol{\xi}-\mathbf{a})
\switch
    \mathcal{F} \left[\mathrm{e}^{\mathrm{i} \mathbf{a}\cdot\boldsymbol{\omega}}f(\mathbf{x})\right](\boldsymbol{\omega})&=f(\boldsymbol{\omega}-\mathbf{a})
    \end{lralign}
\end{proposition}
\begin{proof}
    \ding{172}时移性质，
    \begin{lralign}
    &\mathcal{F} \left[f(\mathbf{x}-\mathbf{a})\right](\boldsymbol{\xi})\\
    &=\int_{\mathbb{R}^n}f(\mathbf{x}-\mathbf{a})\mathrm{e}^{-2\pi \mathrm{i} \mathbf{x}\cdot\boldsymbol{\xi}}\,d\mathbf{x}\\
    =&\int_{\mathbb{R}^n}f(\mathbf{y})\mathrm{e}^{-2\pi \mathrm{i} (\mathbf{y}+\mathbf{a})\cdot\boldsymbol{\xi}}\,d\mathbf{y}(\mathbf{y}=\mathbf{x}-\mathbf{a})\\
    =&\mathrm{e}^{-2\pi \mathrm{i} \mathbf{a}\cdot\boldsymbol{\xi}}\int_{\mathbb{R}^n}f(\mathbf{y})\mathrm{e}^{-2\pi \mathrm{i} \mathbf{y}\cdot\boldsymbol{\xi}}\,d\mathbf{y}\\
    =&\mathrm{e}^{-2\pi \mathrm{i} \mathbf{a}\cdot\boldsymbol{\xi}}\mathcal{F} f(\boldsymbol{\xi})
\switch
    &\mathcal{F} \left[f(\mathbf{x}-\mathbf{a})\right](\boldsymbol{\omega})\\
    =&\int_{\mathbb{R}^n}f(\mathbf{x}-\mathbf{a})\mathrm{e}^{-\mathrm{i} \mathbf{x}\cdot\boldsymbol{\omega}}\,d\mathbf{x}\\
    =&\int_{\mathbb{R}^n}f(\mathbf{y})\mathrm{e}^{-\mathrm{i} (\mathbf{y}+\mathbf{a})\cdot\boldsymbol{\omega}}\,d\mathbf{y}(\mathbf{y}=\mathbf{x}-\mathbf{a})\\
    =&\mathrm{e}^{-\mathrm{i} \mathbf{a}\cdot\boldsymbol{\omega}}\int_{\mathbb{R}^n}f(\mathbf{y})\mathrm{e}^{-\mathrm{i} \mathbf{y}\cdot\boldsymbol{\omega}}\,d\mathbf{y}\\
    =&\mathrm{e}^{-\mathrm{i} \mathbf{a}\cdot\boldsymbol{\omega}}\mathcal{F} f(\boldsymbol{\omega})
    \end{lralign}
    \ding{173}频移性质，
    \begin{lralign}
    &\mathcal{F} \left[\mathrm{e}^{2\pi \mathrm{i} \mathbf{a}\cdot\boldsymbol{\xi}}f(\mathbf{x})\right](\boldsymbol{\xi})\\
    =&\int_{\mathbb{R}^n}\mathrm{e}^{2\pi \mathrm{i} \mathbf{a}\cdot\boldsymbol{\xi}}f(\mathbf{x})\mathrm{e}^{-2\pi \mathrm{i} \mathbf{x}\cdot\boldsymbol{\xi}}\,d\mathbf{x}\\
    =&\int_{\mathbb{R}^n}f(\mathbf{x})\mathrm{e}^{-2\pi \mathrm{i} (\mathbf{x}-\mathbf{a})\cdot\boldsymbol{\xi}}\,d\mathbf{x}\\
    =&f(\boldsymbol{\xi}-\mathbf{a})
\switch
    &\mathcal{F} \left[\mathrm{e}^{\mathrm{i} \mathbf{a}\cdot\boldsymbol{\omega}}f(\mathbf{x})\right](\boldsymbol{\omega})\\
    =&\int_{\mathbb{R}^n}\mathrm{e}^{\mathrm{i} \mathbf{a}\cdot\boldsymbol{\omega}}f(\mathbf{x})\mathrm{e}^{-\mathrm{i} \mathbf{x}\cdot\boldsymbol{\omega}}\,d\mathbf{x}\\
    =&\int_{\mathbb{R}^n}f(\mathbf{x})\mathrm{e}^{-\mathrm{i} (\mathbf{x}-\mathbf{a})\cdot\boldsymbol{\omega}}\,d\mathbf{x}\\
    =&f(\boldsymbol{\omega}-\mathbf{a})
    \end{lralign}
\end{proof}

\begin{proposition}[伸缩定理]
    设$f(\mathbf{x}):\mathbb{R}^n\to\mathbb{C}$，$A$为$n$阶可逆矩阵，则有：
    \begin{lralign}
    \mathcal{F} \left[f(A\mathbf{x})\right](\boldsymbol{\xi})&=\frac{1}{|\det A|}\mathcal{F} f\left(A^{-T}\boldsymbol{\xi}\right)
    \switch
    \mathcal{F} \left[f(A\mathbf{x})\right](\boldsymbol{\omega})&=\frac{1}{|\det A|}\mathcal{F} f\left(A^{-T}\boldsymbol{\omega}\right)
    \end{lralign}
    特别地，当$A$是对角阵时，设$A=\text{diag}(a_1,a_2,\cdots,a_n)$，则有：
    \begin{lralign}
    &\mathcal{F} \left[f(a_1 x_1,a_2 x_2,\cdots,a_n x_n)\right](\boldsymbol{\xi})\\
    =&\frac{1}{|a_1 a_2\cdots a_n|}\mathcal{F} f\left(\frac{\xi_1}{a_1},\frac{\xi_2}{a_2},\cdots,\frac{\xi_n}{a_n}\right)
\switch
    &\mathcal{F} \left[f(a_1 x_1,a_2 x_2,\cdots,a_n x_n)\right](\boldsymbol{\omega})\\
    =&\frac{1}{|a_1 a_2\cdots a_n|}\mathcal{F} f\left(\frac{\omega_1}{a_1},\frac{\omega_2}{a_2},\cdots,\frac{\omega_n}{a_n}\right)
    \end{lralign}
    对于旋转矩阵$R$（即满足$R^T R=I$的矩阵），有：
    \begin{lralign}
    \mathcal{F} \left[f(R\mathbf{x})\right](\boldsymbol{\xi})&=\mathcal{F} f\left(R\boldsymbol{\xi}\right)
\switch
    \mathcal{F} \left[f(R\mathbf{x})\right](\mathbf{\omega})&=\mathcal{F} f\left(R\mathbf{\omega}\right)
    \end{lralign}
    这表明时域旋转，频域也会做相应旋转。
\end{proposition}
\begin{proof}
    \begin{lralign}
    &\mathcal{F} \left[f(A\mathbf{x})\right](\boldsymbol{\xi})\\
    =&\int_{\mathbb{R}^n}f(A\mathbf{x})\mathrm{e}^{-2\pi \mathrm{i} \mathbf{x}\cdot\boldsymbol{\xi}}\,d\mathbf{x}\\
    =&\int_{\mathbb{R}^n}f(\mathbf{y})\mathrm{e}^{-2\pi \mathrm{i} A^{-1}\mathbf{y}\cdot\boldsymbol{\xi}}\,\frac{d\mathbf{y}}{|\det A|}(\mathbf{y}=A\mathbf{x})\\
    =&\frac{1}{|\det A|}\int_{\mathbb{R}^n}f(\mathbf{y})\mathrm{e}^{-2\pi \mathrm{i} \mathbf{y}\cdot A^{-T}\boldsymbol{\xi}}\,d\mathbf{y}\\
    =&\frac{1}{|\det A|}\mathcal{F} f\left(A^{-T}\boldsymbol{\xi}\right)
\switch
    &\mathcal{F} \left[f(A\mathbf{x})\right](\boldsymbol{\omega})\\
    =&\int_{\mathbb{R}^n}f(A\mathbf{x})\mathrm{e}^{-\mathrm{i} \mathbf{x}\cdot\boldsymbol{\omega}}\,d\mathbf{x}\\
    =&\int_{\mathbb{R}^n}f(\mathbf{y})\mathrm{e}^{-\mathrm{i} A^{-1}\mathbf{y}\cdot\boldsymbol{\omega}}\,\frac{d\mathbf{y}}{|\det A|}(\mathbf{y}=A\mathbf{x})\\
    =&\frac{1}{|\det A|}\int_{\mathbb{R}^n}f(\mathbf{y})\mathrm{e}^{-\mathrm{i} \mathbf{y}\cdot A^{-T}\boldsymbol{\omega}}\,d\mathbf{y}\\
    =&\frac{1}{|\det A|}\mathcal{F} f\left(A^{-T}\boldsymbol{\omega}\right)
    \end{lralign}
    请读者自行验证$A$为对角矩阵和旋转矩阵情况下的公式。
\end{proof}

\begin{proposition}[微分性质]
    设$f(\mathbf{x}):\mathbb{R}^n\to\mathbb{C}$一阶连续可导，则有：
    \begin{lralign}
    \mathcal{F} \left(\frac{\partial f}{\partial x_j}\right)(\boldsymbol{\xi})&=2\pi \mathrm{i} \xi_j \mathcal{F} f(\boldsymbol{\xi})
\switch
    \mathcal{F} \left(\frac{\partial f}{\partial x_j}\right)(\boldsymbol{\omega})&=\mathrm{i} \omega_j \mathcal{F} f(\boldsymbol{\omega})
    \end{lralign}
    易见在高阶混合偏导数存在时，有：
    \begin{lralign}
    \mathcal{F} \left(\frac{\partial^{\alpha} f}{\partial x_1^{\alpha_1}\partial x_2^{\alpha_2}\cdots \partial x_n^{\alpha_n}}\right)(\boldsymbol{\xi})\\
    =(2\pi \mathrm{i})^{\alpha} \xi_1^{\alpha_1}\xi_2^{\alpha_2}\cdots \xi_n^{\alpha_n} \mathcal{F} f(\boldsymbol{\xi})
\switch
    \mathcal{F} \left(\frac{\partial^{\alpha} f}{\partial x_1^{\alpha_1}\partial x_2^{\alpha_2}\cdots \partial x_n^{\alpha_n}}\right)(\boldsymbol{\omega})\\
    =\mathrm{i}^{\alpha} \omega_1^{\alpha_1}\omega_2^{\alpha_2}\cdots \omega_n^{\alpha_n} \mathcal{F} f(\boldsymbol{\omega})
    \end{lralign}
    其中$\alpha=\alpha_1+\alpha_2+\cdots +\alpha_n$。
\end{proposition}
\begin{proof}
    以$j=1$为例，其他分量类似。简洁起见，以下不区分重积分时$d\mathbf{x}$的含义。
    {\color{blue}\begin{align*}
    \mathcal{F} \left(\frac{\partial f}{\partial x_1}\right)(\boldsymbol{\xi})=&\int_{\mathbb{R}^n}\frac{\partial f}{\partial x_1}(\mathbf{x})\mathrm{e}^{-2\pi \mathrm{i} \mathbf{x}\cdot\boldsymbol{\xi}}\,d\mathbf{x}\\
    =&\int_{\mathbb{R}^{n-1}}\mathrm{e}^{-2\pi \mathrm{i} (x_2\xi_2+\cdots +x_n\xi_n)}\,d\mathbf{x}\int_{-\infty}^{\infty}\frac{\partial f}{\partial x_1}(\mathbf{x})\mathrm{e}^{-2\pi \mathrm{i} x_1 \xi_1}\,dx_1\\
    =&\int_{\mathbb{R}^{n-1}}\mathrm{e}^{-2\pi \mathrm{i} (x_2\xi_2+\cdots +x_n\xi_n)}\left[f(\mathbf{x})\mathrm{e}^{-2\pi \mathrm{i} x_1 \xi_1}\Big|_{x_1=-\infty}^{x_1=\infty}+2\pi \mathrm{i} \xi_1 \int_{-\infty}^{\infty}f(\mathbf{x})\mathrm{e}^{-2\pi \mathrm{i} x_1 \xi_1}\,dx_1\right]\,d\mathbf{x}\\
    =&2\pi \mathrm{i} \xi_1 \int_{\mathbb{R}^n}f(\mathbf{x})\mathrm{e}^{-2\pi \mathrm{i} \mathbf{x}\cdot\boldsymbol{\xi}}\,d\mathbf{x}\\
    =&2\pi \mathrm{i} \xi_1 \mathcal{F} f(\boldsymbol{\xi})
    \end{align*}}
    {\color{red}
    \begin{align*}
    \mathcal{F} \left(\frac{\partial f}{\partial x_1}\right)(\boldsymbol{\omega})=&\int_{\mathbb{R}^n}\frac{\partial f}{\partial x_1}(\mathbf{x})\mathrm{e}^{-\mathrm{i} \mathbf{x}\cdot\boldsymbol{\omega}}\,d\mathbf{x}\\
    =&\int_{\mathbb{R}^{n-1}}\mathrm{e}^{-\mathrm{i} (x_2\omega_2+\cdots +x_n\omega_n)}\,d\mathbf{x}\int_{-\infty}^{\infty}\frac{\partial f}{\partial x_1}(\mathbf{x})\mathrm{e}^{-\mathrm{i} x_1 \omega_1}\,dx_1\\
    =&\int_{\mathbb{R}^{n-1}}\mathrm{e}^{-\mathrm{i} (x_2\omega_2+\cdots +x_n\omega_n)}\left[f(\mathbf{x})\mathrm{e}^{-\mathrm{i} x_1 \omega_1}\Big|_{x_1=-\infty}^{x_1=\infty}+\mathrm{i} \omega_1 \int_{-\infty}^{\infty}f(\mathbf{x})\mathrm{e}^{-\mathrm{i} x_1 \omega_1}\,dx_1\right]\,d\mathbf{x}\\
    =&\mathrm{i} \omega_1 \int_{\mathbb{R}^n}f(\mathbf{x})\mathrm{e}^{-\mathrm{i} \mathbf{x}\cdot\boldsymbol{\omega}}\,d\mathbf{x}\\
    =&\mathrm{i} \omega_1 \mathcal{F} f(\boldsymbol{\omega})
    \end{align*}}
\end{proof}

\begin{definition}[场论中的算子]
    $\nabla$为哈密顿算子（或称nabla算子），$\Delta=\nabla\cdot \nabla$为拉普拉斯算子，
    \begin{align*}
        \nabla=\left(\frac{\partial }{\partial x_1},\frac{\partial }{\partial x_2},\cdots,\frac{\partial }{\partial x_n}\right),
        \Delta=\frac{\partial^2 }{\partial x_1^2}+\frac{\partial^2 }{\partial x_2^2}+\cdots+\frac{\partial^2 }{\partial x_n^2}
    \end{align*}
    场论中，梯度、散度、旋度等运算都可以用$\nabla$表示：
    \begin{itemize}
        \item 梯度：$\nabla f=\text{grad}f$
        \item 散度：$\nabla\cdot \mathbf{F}=\text{div}\mathbf{F}$，其中$\mathbf{F}=(F_1,F_2,\cdots,F_n)$
        \item 旋度（仅在$\mathbb{R}^3$中定义）：$\nabla\times \mathbf{F}=\text{curl}\mathbf{F}$（一些文献中也将旋度算子记为$\text{rot}$）
    \end{itemize}
\end{definition}

\begin{corollary}[场论中的微分性质]
    由微分性质，我们依次对比向量值函数的各个分量。对于梯度$\nabla f=\text{grad}f$，有：
    \begin{lralign}
    \mathcal{F} \left(\nabla f\right)(\boldsymbol{\xi})&=2\pi \mathrm{i} \boldsymbol{\xi} \mathcal{F} f(\boldsymbol{\xi})
\switch
    \mathcal{F} \left(\nabla f\right)(\boldsymbol{\omega})&=\mathrm{i} \boldsymbol{\omega} \mathcal{F} f(\boldsymbol{\omega})
    \end{lralign}
    对于散度$\nabla\cdot \mathbf{F}=\text{div}\mathbf{F}$，其中$\mathbf{F}=(F_1,F_2,\cdots,F_n)$，有：
    \begin{lralign}
    \mathcal{F} \left(\nabla\cdot \mathbf{F}\right)(\boldsymbol{\xi})&=2\pi \mathrm{i} \boldsymbol{\xi}\cdot \mathcal{F} \mathbf{F}(\boldsymbol{\xi})
\switch
    \mathcal{F} \left(\nabla\cdot \mathbf{F}\right)(\boldsymbol{\omega})&=\mathrm{i} \boldsymbol{\omega}\cdot \mathcal{F} \mathbf{F}(\boldsymbol{\omega})
    \end{lralign}
    在$\mathbb{R}^3$中，对于旋度$\nabla\times \mathbf{F}=\text{curl}\mathbf{F}$，有：
    \begin{lralign}
    \mathcal{F} \left(\nabla\times \mathbf{F}\right)(\boldsymbol{\xi})&=2\pi \mathrm{i} \boldsymbol{\xi}\times \mathcal{F} \mathbf{F}(\boldsymbol{\xi})
\switch
    \mathcal{F} \left(\nabla\times \mathbf{F}\right)(\boldsymbol{\omega})&=\mathrm{i} \boldsymbol{\omega}\times \mathcal{F} \mathbf{F}(\boldsymbol{\omega})
    \end{lralign}
    对于拉普拉斯算子$\Delta f$，有：
    \begin{lralign}
    \mathcal{F} \left(\Delta f\right)(\boldsymbol{\xi})&=-(2\pi)^2 |\boldsymbol{\xi}|^2 \mathcal{F} f(\boldsymbol{\xi})
\switch
    \mathcal{F} \left(\Delta f\right)(\boldsymbol{\omega})&=-|\boldsymbol{\omega}|^2 \mathcal{F} f(\boldsymbol{\omega})
    \end{lralign}
\end{corollary}

我们还可以定义n维空间$\mathbb{R}^n$上的\textbf{卷积}。与一维的情况类似，由于卷积的存在性可以由多种充分条件保证，我们直接从形式上给出卷积的定义，而不讨论$f,g$具体的范围。
\begin{definition}[卷积]
    设$f(\mathbf{x}),g(\mathbf{x}):\mathbb{R}^n\to\mathbb{C}$，它们的卷积定义为：
    \begin{align*}
        (f*g)(\mathbf{x})=\int_{\mathbb{R}^n}f(\mathbf{y})g(\mathbf{x}-\mathbf{y})\,d\mathbf{y}
    \end{align*}
\end{definition}

\begin{theorem}[卷积定理]
    设$f(\mathbf{x}),g(\mathbf{x}):\mathbb{R}^n\to\mathbb{C}$，则有：
    \begin{lralign}
    &\mathcal{F} (f*g)(\boldsymbol{\xi})=\mathcal{F} f(\boldsymbol{\xi})\cdot \mathcal{F} g(\boldsymbol{\xi})\\
    &\mathcal{F} f(\boldsymbol{\xi})\mathcal{F} g(\boldsymbol{\xi})=\mathcal{F} (f*g)(\boldsymbol{\xi})
\switch
    &\mathcal{F} (f*g)(\boldsymbol{\omega})=\mathcal{F} f(\boldsymbol{\omega})\cdot \mathcal{F} g(\boldsymbol{\omega})\\
    &\mathcal{F} f(\boldsymbol{\omega})\mathcal{F} g(\boldsymbol{\omega})=\frac{1}{(2\pi)^n}\mathcal{F} (f*g)(\boldsymbol{\omega})
    \end{lralign}
\end{theorem}
\begin{proof}
    \ding{172}时域卷积：
\begin{lralign}
    &\mathcal{F} (f*g)(\boldsymbol{\xi})\\
    =&\int_{\mathbb{R}^n}(f*g)(\mathbf{x})\mathrm{e}^{-2\pi \mathrm{i} \mathbf{x}\cdot\boldsymbol{\xi}}\,d\mathbf{x}\\
    =&\int_{\mathbb{R}^n}\mathrm{e}^{-2\pi \mathrm{i} \mathbf{x}\cdot\boldsymbol{\xi}}\,d\mathbf{x}\int_{\mathbb{R}^n}f(\mathbf{y})g(\mathbf{x}-\mathbf{y})\,d\mathbf{y}\\
    =&\int_{\mathbb{R}^n}f(\mathbf{y})\,d\mathbf{y}\int_{\mathbb{R}^n}g(\mathbf{x}-\mathbf{y})\mathrm{e}^{-2\pi \mathrm{i} \mathbf{x}\cdot\boldsymbol{\xi}}\,d\mathbf{x}\\
    =&\int_{\mathbb{R}^n}f(\mathbf{y})\,d\mathbf{y}\int_{\mathbb{R}^n}g(\mathbf{z})\mathrm{e}^{-2\pi \mathrm{i} (\mathbf{z}+\mathbf{y})\cdot\boldsymbol{\xi}}\,d\mathbf{z}\\
    =&\int_{\mathbb{R}^n}f(\mathbf{y})\mathrm{e}^{-2\pi \mathrm{i} \mathbf{y}\cdot\boldsymbol{\xi}}\,d\mathbf{y}\int_{\mathbb{R}^n}g(\mathbf{z})\mathrm{e}^{-2\pi \mathrm{i} \mathbf{z}\cdot\boldsymbol{\xi}}\,d\mathbf{z}\\
    =&\mathcal{F} f(\boldsymbol{\xi})\cdot \mathcal{F} g(\boldsymbol{\xi})
    \switch
    &\mathcal{F} (f*g)(\boldsymbol{\omega})\\
    =&\int_{\mathbb{R}^n}(f*g)(\mathbf{x})\mathrm{e}^{-\mathrm{i} \mathbf{x}\cdot\boldsymbol{\omega}}\,d\mathbf{x}\\
    =&\int_{\mathbb{R}^n}f(\mathbf{y})\,d\mathbf{y}\int_{\mathbb{R}^n}g(\mathbf{x}-\mathbf{y})\mathrm{e}^{-\mathrm{i} \mathbf{x}\cdot\boldsymbol{\omega}}\,d\mathbf{x}\\
    =&\int_{\mathbb{R}^n}f(\mathbf{y})\,d\mathbf{y}\int_{\mathbb{R}^n}g(\mathbf{x}-\mathbf{y})\mathrm{e}^{-\mathrm{i} \mathbf{x}\cdot\boldsymbol{\omega}}\,d\mathbf{x}\\
    =&\int_{\mathbb{R}^n}f(\mathbf{y})\,d\mathbf{y}\int_{\mathbb{R}^n}g(\mathbf{z})\mathrm{e}^{-\mathrm{i} (\mathbf{z}+\mathbf{y})\cdot\boldsymbol{\omega}}\,d\mathbf{z}\\
    =&\int_{\mathbb{R}^n}f(\mathbf{y})\mathrm{e}^{-\mathrm{i} \mathbf{y}\cdot\boldsymbol{\omega}}\,d\mathbf{y}\int_{\mathbb{R}^n}g(\mathbf{z})\mathrm{e}^{-\mathrm{i} \mathbf{z}\cdot\boldsymbol{\omega}}\,d\mathbf{z}\\
    =&\mathcal{F} f(\boldsymbol{\omega})\cdot \mathcal{F} g(\boldsymbol{\omega})
    \end{lralign}
\ding{173}频域卷积：\\
仿照\ref{sub:properties_of_convolution}中的证明。根据时域卷积定理$\mathcal{F} (f*g)=\mathcal{F} f\mathcal{F} g$，
令$f=\mathcal{F} \mathfrak{f},g=\mathcal{F} \mathfrak{g}$，则
\begin{lralign}
    \mathcal{F} \mathfrak{f}*\mathcal{F} \mathfrak{g}&=\mathcal{F} ^{-1}\left[\mathcal{F}\mathcal{F} \mathfrak{f}\cdot\mathcal{F} \mathcal{F} \mathfrak{g}\right]\\
    &=\mathcal{F} ^{-1}\left(\mathfrak{f}^- \mathfrak{g}^- \right)\\
    &=\mathcal{F} (\mathfrak{fg})
\switch
    \mathcal{F} \mathfrak{f}*\mathcal{F} \mathfrak{g}&=\mathcal{F} ^{-1}\left[\mathcal{F}\mathcal{F} \mathfrak{f}\cdot\mathcal{F} \mathcal{F} \mathfrak{g}\right]\\
    &=\mathcal{F} ^{-1}\left((2\pi)^{2n} \mathfrak{f}^- \mathfrak{g}^- \right)\\
    &= (2\pi)^n\mathcal{F} (\mathfrak{fg})
\end{lralign}
\end{proof}

\begin{theorem}[普朗歇尔定理]
    设$f(\mathbf{x}),\in L^1(\mathbb{R}^n)\cap L^2(\mathbb{R}^n)$，则有：
    \begin{lralign}
    \int_{\mathbb{R}^n}|f(\mathbf{x})|^2\,d\mathbf{x}&=\int_{\mathbb{R}^n}|\mathcal{F} f(\boldsymbol{\xi})|^2\,d\boldsymbol{\xi}
\switch
    \int_{\mathbb{R}^n}|f(\mathbf{x})|^2\,d\mathbf{x}&=\frac{1}{(2\pi)^n}\int_{\mathbb{R}^n}|\mathcal{F} f(\boldsymbol{\omega})|^2\,d\boldsymbol{\omega}
    \end{lralign}
\end{theorem}
\begin{proof}
    与一维的情况类似，我们给出更一般的结论，它实际上指出空间域与变换域上的内积的数量关系。设$f,g\in L^1(\mathbb{R}^n)\cap L^2(\mathbb{R}^n)$，则有：
    \begin{lralign}
    &\int_{\mathbb{R}^n}\mathcal{F} f(\boldsymbol{\xi})\overline{\mathcal{F} g(\boldsymbol{\xi})}\,d\mathbf{x}\\
    =&\int_{\mathbb{R}^n}\,d\boldsymbol{\xi}\int_{\mathbb{R}^n}f(\mathbf{x})\mathrm{e}^{-2\pi \mathrm{i} \mathbf{x}\cdot\boldsymbol{\xi}}\,d\mathbf{x}\int_{\mathbb{R}^n}\overline{g(\mathbf{y})}\mathrm{e}^{2\pi \mathrm{i} \mathbf{y}\cdot\boldsymbol{\xi}}\,d\mathbf{y}\\
    =&\int_{\mathbb{R}^n}f(\mathbf{x})\,d\mathbf{x}\int_{\mathbb{R}^n}\overline{g(\mathbf{y})}\,d\mathbf{y}\int_{\mathbb{R}^n}\mathrm{e}^{-2\pi \mathrm{i} (\mathbf{x}-\mathbf{y})\cdot\boldsymbol{\xi}}\,d\boldsymbol{\xi}\\
    =&\int_{\mathbb{R}^n}f(\mathbf{x})\,d\mathbf{x}\int_{\mathbb{R}^n}\overline{g(\mathbf{y})}\delta(\mathbf{x}-\mathbf{y})\,d\mathbf{y}\\
    =&\int_{\mathbb{R}^n}f(\mathbf{x})\overline{g(\mathbf{x})}\,d\mathbf{x}
\switch
    &\int_{\mathbb{R}^n}\mathcal{F} f(\boldsymbol{\omega})\overline{\mathcal{F} g(\boldsymbol{\omega})}\,d\mathbf{x}\\
    =&\int_{\mathbb{R}^n}\,d\boldsymbol{\omega}\int_{\mathbb{R}^n}f(\mathbf{x})\mathrm{e}^{-\mathrm{i} \mathbf{x}\cdot\boldsymbol{\omega}}\,d\mathbf{x}\int_{\mathbb{R}^n}\overline{g(\mathbf{y})}\mathrm{e}^{\mathrm{i} \mathbf{y}\cdot\boldsymbol{\omega}}\,d\mathbf{y}\\
    =&\int_{\mathbb{R}^n}f(\mathbf{x})\,d\mathbf{x}\int_{\mathbb{R}^n}\overline{g(\mathbf{y})}\,d\mathbf{y}\int_{\mathbb{R}^n}\mathrm{e}^{-\mathrm{i} (\mathbf{x}-\mathbf{y})\cdot\boldsymbol{\omega}}\,d\boldsymbol{\omega}\\
    =&\int_{\mathbb{R}^n}f(\mathbf{x})\,d\mathbf{x}\int_{\mathbb{R}^n}\overline{g(\mathbf{y})}(2\pi)^n \delta(\mathbf{x}-\mathbf{y})\,d\mathbf{y}\\
    =&(2\pi)^n \int_{\mathbb{R}^n}f(\mathbf{x})\overline{g(\mathbf{x})}\,d\mathbf{x}
    \end{lralign}
    取$g=f$即得所需结论。
\end{proof}

\subsection{多维施瓦兹函数与多维分布}\label{sub:multi_schwartz}

首先，我们将\textbf{施瓦兹函数}的定义推广到多维空间$\mathbb{R}^n$中。
\begin{definition}[多维施瓦兹函数]
    设$\varphi(\mathbf{x}):\mathbb{R}^n\to\mathbb{C}$，如果\begin{itemize}
        \item $\varphi(\mathbf{x})\in C^{\infty}(\mathbb{R}^n)$；
        \item 对任意多重指标$\mathcal{J}=(j_1,j_2,\ldots,j_n)$，都有
        \[\lim_{|\mathbf{x}|\to\infty}|\mathbf{x}^{\mathcal{J}} \partial_{\mathcal{J}} \varphi(\mathbf{x})|=0\]
    \end{itemize}
    其中$$\partial_{\mathcal{J}}=\frac{\partial^{j_1+j_2+\cdots +j_n}}{\partial x_1^{j_1}\partial x_2^{j_2}\cdots \partial x_n^{j_n}}$$
    则称$\varphi(\mathbf{x})$为\textbf{多维施瓦兹函数}，所有多维施瓦兹函数构成的空间记为$\mathcal{S}(\mathbb{R}^n)$，在不引起歧义的前提下，也可直接记作$\mathcal{S}$。
\end{definition}

$\mathcal{S} $是线性空间，因此可以定义其对偶空间，即\textbf{多维缓增分布空间}，$\mathcal{T} $中的元素称为\textbf{多维缓增分布}：
\begin{definition}[多维缓增分布]
    泛函$T:\mathcal{S} \to \mathbb{C}$为线性映射，将施瓦兹函数$\varphi\in \mathcal{S} $映射为复数$T(\varphi)\in \mathbb{C}$。对于正则分布$T_f$，有
    $$T_f(\varphi)=\langle T_f,\varphi\rangle=\int_{\mathbb{R}^n}f(\mathbf{x})\varphi(\mathbf{x})\,d\mathbf{x}$$
    将缓增分布组成的线性空间记作$\mathcal{T} (\mathbb{R}^n)$，在不引起歧义的前提下，也可直接记作$\mathcal{T} $。
\end{definition}

容易验证$\mathcal{S} $上的卷积总是有定义的；用类似于\ref{sub:fourier_close}中的方法，可以证明$\mathcal{S} $在傅里叶变换下是封闭的，我们仍然定义$\langle \mathcal{F} T,\varphi\rangle=\langle T,\mathcal{F} \varphi\rangle$从而$\mathcal{T} $在傅里叶变换下也是封闭的。典型的施瓦兹函数例如高斯函数$f(\mathbf{x})=\mathrm{e}^{-\pi |\mathbf{x}|^2}$。

我们不再一一讨论作用于分布的各算子，因为它们与一维的情况没有多少差异。对于分布而言，平移与伸缩需要用算子来表示，只是这时平移算子$\tau_{\mathbf{b}}$（其中$\mathbf{b}$为向量）定义为：
\begin{align*}
    (\tau_{\mathbf{b}} f)(\mathbf{x})=f(\mathbf{x}-\mathbf{b})
\end{align*}
伸缩算子$\sigma_A$（其中$A$为矩阵）定义为：
\begin{align*}
    (\sigma_A f)(\mathbf{x})=f(A\mathbf{x})
\end{align*}

下面以这两个算子为例，给出它们作用于缓增分布$T\in\mathcal{T} $的定义，从而看出它们与一维情况下的相似性。
\begin{example}[平移算子]
    设$f,\varphi\in\mathcal{S} ,T_f\in\mathcal{T} $，自然可以定义$\tau_{\mathbf{b}} T_f=T_{\tau_{\mathbf{b}} f}$，
    \begin{align*}
        \langle \tau_{\mathbf{b}} T_f,\varphi\rangle=&\int_{\mathbb{R}^n}f(\mathbf{x}-\mathbf{b})\varphi(\mathbf{x})\,d\mathbf{x}\\
        =&\int_{\mathbb{R}^n}f(\mathbf{y})\varphi(\mathbf{y}+\mathbf{b})\,d\mathbf{y}\quad (\mathbf{y}=\mathbf{x}-\mathbf{b})\\
        =&\langle T_f,\tau_{-\mathbf{b}}\varphi\rangle
    \end{align*}
    因此我们定义：
    \[\langle \tau_{\mathbf{b}} T,\varphi\rangle=\langle T,\tau_{-\mathbf{b}}\varphi\rangle,T\in\mathcal{T} ,\varphi\in\mathcal{S} \]
\end{example}

\begin{example}[伸缩算子]
    设$f,\varphi\in\mathcal{S} ,T_f\in\mathcal{T} $，自然可以定义$\sigma_A T_f=T_{\sigma_A f}$，
    \begin{align*}
        \langle \sigma_A T_f,\varphi\rangle=&\int_{\mathbb{R}^n}f(A\mathbf{x})\varphi(\mathbf{x})\,d\mathbf{x}\\
        =&\int_{\mathbb{R}^n}f(\mathbf{y})\varphi(A^{-1}\mathbf{y})\,\frac{d\mathbf{y}}{|\det A|}\quad (\mathbf{y}=A\mathbf{x})\\
        =&\frac{1}{|\det A|}\langle T_f,\sigma_{A^{-1}}\varphi\rangle
    \end{align*}
    因此我们定义：
    \[\langle \sigma_A T,\varphi\rangle=\frac{1}{|\det A|}\langle T,\sigma_{A^{-1}}\varphi\rangle,T\in\mathcal{T} ,\varphi\in\mathcal{S} \]
\end{example}
在一维情况下，我们有：

\[\langle \tau_a T,\varphi\rangle=\langle T,\tau_{-a}\varphi\rangle,\langle\sigma_{a} T,\varphi\rangle=\frac{1}{|a|}\langle T,\sigma_{1/a}\varphi\rangle,T\in\mathcal{T} ,\varphi\in\mathcal{S} \]
可以看到，如果将标量$b$推广为向量$\mathbf{b}$，而将线性变换（的矩阵表示）由一维矩阵$a$推广为$n$维矩阵$A$，则一维情况下的平移和伸缩算子的定义就是多维空间中的特例。

借助上一小节中对函数的傅里叶变换的讨论，读者可以自行构建出分布的傅里叶变换的各种性质。

下面，我们来讨论一些典型的分布。首先是多维狄拉克分布：
\begin{definition}[多维狄拉克分布]
    设$\mathbf{x}_0\in\mathbb{R}^n$，则多维狄拉克分布$\delta_{\mathbf{x}_0}\in\mathcal{T} $定义为：
    \[\langle \delta_{\mathbf{x}_0},\varphi\rangle=\varphi(\mathbf{x}_0),\forall \varphi\in\mathcal{S} \]
    特别地，$\delta_{\mathbf{0}}$记为$\delta$，称为\textbf{多维狄拉克分布}。

    形式上，可以将多维狄拉克分布表示为：
    \[\delta_{\mathbf{x}_0}(\mathbf{x})=\begin{cases}
    \infty,&\mathbf{x}=\mathbf{x}_0\\
    0,&\mathbf{x}\neq \mathbf{x}_0
    \end{cases},\int_{\mathbb{R}^n}\delta(\mathbf{x})\,d\mathbf{x}=1\]
    即$\delta_{\mathbf{x}_0}(\mathbf{x})=\tau_{\mathbf{x}_0}\delta(\mathbf{x})$。
\end{definition}
\begin{proposition}[多维狄拉克分布的性质]
    多维狄拉克分布具有以下性质：
    \begin{itemize}
        \item 多维狄拉克分布仍为卷积幺元，即$f*\delta=\delta*f=f$；
        \item 多维狄拉克分布是可分的，即
        \[\delta(\mathbf{x})=\delta(x_1,x_2,\ldots,x_n)=\delta(x_1)\delta(x_2)\cdots \delta(x_n)\]
        \item 多维狄拉克分布的伸缩性质：
        \[\sigma_A \delta=\frac{1}{|\det A|}\delta\]
        特别地，\[\delta(a_1 x_1,a_2 x_2,\ldots,a_n x_n)=\frac{1}{|a_1 a_2 \cdots a_n|}\delta(x_1,x_2,\ldots,x_n)\]
        并且$\delta$是偶分布，即$\delta^-=\delta$；
        \item 多维狄拉克分布的傅里叶变换：
        \begin{lralign}
        \mathcal{F} \delta(\boldsymbol{\xi})&=\mathds{1}
\switch
        \mathcal{F} \delta(\boldsymbol{\omega})&=\mathds{1}
        \end{lralign}
    \end{itemize}
\end{proposition}
\begin{proof}
    分布与函数的卷积定义为：
    \[\langle f*T,\varphi\rangle=\langle T,(f^-)*\varphi\rangle,f,\varphi\in\mathcal{S} ,T\in\mathcal{T} \]
    并且这种卷积总满足交换律，即$f*T=T*f$。因此有：
    \begin{align*}
        \langle f*\delta,\varphi\rangle=\langle \delta,(f^-)*\varphi\rangle=((f^-)*\varphi)(\mathbf{0})=\int_{\mathbb{R}^n}f(\mathbf{y})\varphi(\mathbf{y})\,d\mathbf{y}=\langle f,\varphi\rangle
    \end{align*}
    这就证明了$f*\delta=\delta*f=f$。

    类似地，有：
    \begin{align*}
        \langle \delta(x_1)\delta(x_2)\cdots \delta(x_n),\varphi(\mathbf{x})\rangle&=\int_{\mathbb{R}^n}\delta(x_1)\delta(x_2)\cdots \delta(x_n)\varphi(\mathbf{x})\,d\mathbf{x}\\
        &=\int_{\mathbb{R}^{n-1}}\delta(x_2)\cdots \delta(x_n)\,dx_2\cdots dx_n\int_{-\infty}^{\infty}\delta(x_1)\varphi(\mathbf{x})\,dx_1\\
        &=\int_{\mathbb{R}^{n-1}}\delta(x_2)\cdots \delta(x_n)\varphi(0,x_2,\ldots,x_n)\,dx_2\cdots dx_n\\
        &\vdots\\
        &=\varphi(0,0,\ldots,0)\\
        &=\langle \delta(\mathbf{x}),\varphi(\mathbf{x})\rangle
    \end{align*}
    即$\delta(\mathbf{x})=\delta(x_1)\delta(x_2)\cdots \delta(x_n)$。

    对于伸缩性质，直接套用作用于分布的伸缩算子的定义：
    \begin{align*}
        \langle \sigma_A \delta,\varphi\rangle=&\frac{1}{|\det A|}\langle \delta,\sigma_{A^{-1}}\varphi\rangle\\
        =&\frac{1}{|\det A|}(\sigma_{A^{-1}}\varphi)(\mathbf{0})\\
        =&\frac{1}{|\det A|}\varphi(A^{-1}\mathbf{0})\\
        =&\frac{1}{|\det A|}\varphi(\mathbf{0})\\
        =&\left\langle \frac{1}{|\det A|}\delta,\varphi\right\rangle
    \end{align*}

    最后，求多维狄拉克分布的傅里叶变换：
    \begin{lralign}
    \langle \mathcal{F} \delta,\varphi\rangle&=\langle \delta,\mathcal{F} \varphi\rangle\\
    &=(\mathcal{F} \varphi)(\mathbf{0})\\
    &=\int_{\mathbb{R}^n}\varphi(\mathbf{x})\,d\mathbf{x}\\
    &=\langle \mathds{1},\varphi\rangle
\switch
    \langle \mathcal{F} \delta,\varphi\rangle&=\langle \delta,\mathcal{F} \varphi\rangle\\
    &=(\mathcal{F} \varphi)(\mathbf{0})\\
    &=\frac{1}{(2\pi)^n}\int_{\mathbb{R}^n}\varphi(\mathbf{x})\,d\mathbf{x}\\
    &=\langle \mathds{1},\varphi\rangle
    \end{lralign}
因此多维狄拉克分布的傅里叶变换是常函数1。
\end{proof}

\begin{corollary}
    利用傅里叶变换的对偶性，可以得到：
    \begin{lralign}
    \mathcal{F} \mathds{1}(\boldsymbol{\xi})&=\delta(\boldsymbol{\xi})
\switch
    \mathcal{F} \mathds{1}(\boldsymbol{\omega})&=(2\pi)^n \delta(\boldsymbol{\omega})
    \end{lralign}

    再利用傅里叶变换的频移性质，得到：
    \begin{lralign}
        \mathcal{F} [\mathrm{e}^{2\pi \mathrm{i} \mathbf{x}_0\cdot\mathbf{x}}](\boldsymbol{\xi})&=\delta_{\mathbf{x_0}}
\switch
        \mathcal{F} [\mathrm{e}^{\mathrm{i} \mathbf{x}_0\cdot\mathbf{x}}](\boldsymbol{\omega})&=(2\pi)^n \delta_{\mathbf{x_0}}
    \end{lralign}
    利用欧拉公式求得三角函数的傅里叶变换：
    \begin{lralign}
        \mathcal{F} [\cos(2\pi \mathbf{x}_0\cdot\mathbf{x})](\boldsymbol{\xi})&=\frac{1}{2}[\delta_{\mathbf{x_0}}+\delta_{-\mathbf{x_0}}]
\switch
        \mathcal{F} [\cos(\mathbf{x}_0\cdot\mathbf{x})](\boldsymbol{\omega})&=\frac{(2\pi)^n}{2}[\delta_{\mathbf{x_0}}+\delta_{-\mathbf{x_0}}]
    \end{lralign}
    以及
    \begin{lralign}
        \mathcal{F} [\sin(2\pi \mathbf{x}_0\cdot\mathbf{x})](\boldsymbol{\xi})&=\frac{1}{2\mathrm{i}}[\delta_{\mathbf{x_0}}-\delta_{-\mathbf{x_0}}]
\switch
        \mathcal{F} [\sin(\mathbf{x}_0\cdot\mathbf{x})](\boldsymbol{\omega})&=\frac{(2\pi)^n}{2\mathrm{i}}[\delta_{\mathbf{x_0}}-\delta_{-\mathbf{x_0}}]
    \end{lralign}
\end{corollary}

\subsection{格与二维狄拉克梳状分布}\label{sub:lattice}

为了直观性，我们再这一小节中将仅考虑$\mathbb{R}^2$中的情况。我们首先给出\textbf{格}的定义，讨论它的一些性质，再介绍它在信号与系统中的应用。
\begin{definition}[格]
    任取$\mathbb{R}^2$，的一组基$\{\mathbf{u}_1,\mathbf{u}_2\}$，称形如
    \[\mathbf{p}=p_1\mathbf{u}_1+p_2\mathbf{u}_2,p_1,p_2\in\mathbb{Z}\]
    的点为\textbf{格点}(lattice point)，所有格点组成的集合称为由基$\{\mathbf{u}_1,\mathbf{u}_2\}$生成的\textbf{格}(lattice)，记为$\mathcal{L}$，$\mathcal{L} $构成$\mathbb{Z}$模。

    以$\mathbf{u}_1,\mathbf{u}_2$为邻边的平行四边形称为该格的\textbf{基本单元}或\textbf{基本平行四边形}(fundamental parallelogram)，基本单元的面积为$|\det(\mathbf{u}_1,\mathbf{u}_2)|$，定义为该格的\textbf{面积}（体积），记为$\mathrm{Area}(\mathcal{L} )$。

    记$A=(\mathbf{u}_1,\mathbf{u}_2)$为由基向量组成的可逆矩阵，则
    \[\mathbf{u}_1=A \mathbf{e}_1,\mathbf{u}_2=A \mathbf{e}_2\]
    其中$\mathbf{e}_1,\mathbf{e}_2$为标准正交基。作为列向量，格点可以表示为
    \[\mathbf{p}=A\mathbf{n},\mathbf{n}\in\mathbb{Z}^2\]
    因此，我们也将格$\mathcal{L} $记为$\mathcal{L} =A(\mathbb{Z}^2)$，$\mathrm{Area}(\mathcal{L} )=|\det A|$。
\end{definition}

\begin{figure}[H]
    \centering
    \includegraphics[width=0.8\textwidth]{n_comb}
    \caption{二维的格和狄拉克梳状分布示意图}
\end{figure}

\begin{definition}[狄拉克梳状分布]
    设$\mathcal{L}=A(\mathbb{Z}^2)$为$\mathbb{R}^2$中的格，则由该格生成的\textbf{二维狄拉克梳状分布}(Dirac comb distribution)定义为
    \[\shah_{\mathcal{L} }=\sum_{\mathbf{p}\in\mathcal{L} }\delta_{\mathbf{p}}\]
    特别地，当$\mathcal{L} =\mathbb{Z}^2$时，我们直接记为$\shah$。$\shah_{\mathcal{L} }=\dfrac{1}{|\det A|}\sigma_{A^{-1}}\shah$。

    狄拉克梳状分布看起来像一个钉床，因此又称为\textbf{钉床分布}(bed of nails)。
\end{definition}
\begin{remark}
    $\shah_{\mathcal{L} }=\dfrac{1}{|\det A|}\sigma_{A^{-1}}\shah$的证明如下：
    \begin{align*}
        \frac{1}{|\det A|}\langle \sigma_{A^{-1}}\shah,\varphi\rangle=\frac{1}{|\det A^{-1}|}\langle \shah,\sigma_A \varphi\rangle
        =\sum_{\mathbf{n}\in\mathbb{Z}^2}(\sigma_A \varphi)(\mathbf{n})
        =\sum_{\mathbf{n}\in\mathbb{Z}^2}\varphi(A\mathbf{n})
        =\sum_{\mathbf{p}\in\mathcal{L} }\varphi(\mathbf{p})
        =\langle \shah_{\mathcal{L} },\varphi\rangle
    \end{align*}
\end{remark}

我们将再一次使用狄拉克梳状分布来分析采样信号，只不过这时信号是二维的（例如图像信号），而采样是在格点上进行的。在本小节中，我们仍然使用频率形式的傅里叶变换。

首先，我们给出二维情况下的\textbf{傅里叶级数}以及\textbf{泊松求和公式}：
\begin{definition}[二维傅里叶级数]
    设$f(x_1,x_2)$为$\mathbb{R}^2$中周期为$(T_1,T_2)$的周期函数，即
    \[f(x_1,x_2)=f(x_1+T_1,x_2)=f(x_1,x_2+T_2),\forall (x_1,x_2)\in\mathbb{R}^2\]
    设$f\in L^2([0,T_1]\times[0,T_2])$，则$f$的\textbf{二维傅里叶级数}为
    \[f(x_1,x_2)=\sum_{k_1,k_2=-\infty}^{\infty}c_{k_1,k_2}\mathrm{e}^{2\pi \mathrm{i}\left(\frac{1}{T_1}k_1 x_1+\frac{1}{T_2}k_2 x_2\right)}\]
    其中
    \[c_{k_1,k_2}=\frac{1}{T_1 T_2}\int_0^{T_1}\int_0^{T_2}f(x_1,x_2)\mathrm{e}^{-2\pi \mathrm{i}\left(\frac{1}{T_1}k_1 x_1+\frac{1}{T_2}k_2 x_2\right)}\,dx_1 dx_2\]
    特别地，设$T_1=T_2=1$，则有
    \[f(x_1,x_2)=\sum_{k_1,k_2=-\infty}^{\infty}c_{k_1,k_2}\mathrm{e}^{2\pi \mathrm{i}(k_1 x_1+k_2 x_2)}=\sum_{\mathbf{k}\in\mathbb{Z}^2}c_{\mathbf{k}}\mathrm{e}^{2\pi \mathrm{i}\mathbf{k\cdot x}}\]
    其中
    \[c_{k_1,k_2}=\int_{[0,1]^2}f(x_1,x_2)\mathrm{e}^{-2\pi \mathrm{i}(k_1 x_1+k_2 x_2)}\,dx_1 dx_2=\int_{[0,1]^2}f(\mathbf{x})\mathrm{e}^{-2\pi \mathrm{i} \mathbf{k}\cdot \mathbf{x}}\,d\mathbf{x}\]
\end{definition}
\begin{theorem}[二维泊松求和公式]
    \[\sum_{\mathbf{k}\in\mathbb{Z}^2}\mathcal{F} \varphi(\mathbf{k})=\sum_{\mathbf{k}\in\mathbb{Z}^2}\varphi(\mathbf{k}),\varphi\in\mathcal{S} \]
\end{theorem}
\begin{proof}
    仿照一维情况下的证明（见\ref{sub:poisson_summation}），令$\Phi(\mathbf{x})=(\shah *\varphi)(\mathbf{x})=\sum_{\mathbf{k}\in\mathbb{Z}^2}\varphi(\mathbf{x}-\mathbf{k})$，则$\Phi$为周期为1的二维函数，可以展开为傅里叶级数：
    \[\Phi(\mathbf{x})=\sum_{\mathbf{k}\in\mathbb{Z}^2}\hat{\Phi}(\mathbf{k})\mathrm{e}^{2\pi \mathrm{i} \mathbf{k}\cdot \mathbf{x}}\]
    其中\begin{align*}
        \hat{\Phi}(\mathbf{k})&=\int_{[0,1]^2}\Phi(\mathbf{x})\mathrm{e}^{-2\pi \mathrm{i}\mathbf{k}\cdot \mathbf{x}}\,d\mathbf{x}\\
        &=\int_{[0,1]^2}\sum_{\mathbf{m}\in\mathbb{Z}^2}\varphi(\mathbf{x}-\mathbf{m})\mathrm{e}^{-2\pi \mathrm{i}\mathbf{k}\cdot \mathbf{x}}\,d\mathbf{x}\\
        &=\int_{\mathbb{R}^2}\varphi(\mathbf{y})\mathrm{e}^{-2\pi \mathrm{i}\mathbf{k}\cdot \mathbf{y}}\,d\mathbf{y}\quad (\mathbf{y}=\mathbf{x}-\mathbf{m})\\
        &=\mathcal{F} \varphi(\mathbf{k})
    \end{align*}
    因此，有
    \[\Phi(\mathbf{x})=\sum_{\mathbf{k}\in\mathbb{Z}^2}\mathcal{F} \varphi(\mathbf{k})\mathrm{e}^{2\pi \mathrm{i} \mathbf{k}\cdot \mathbf{x}}\]
    令$\mathbf{x}=\mathbf{0}$，即得所需结论。
\end{proof}

通过二维泊松求和公式，可以得到狄拉克梳状分布的傅里叶变换：
\begin{proposition}[二维狄拉克梳状分布的傅里叶变换]
    \[\mathcal{F} \shah=\shah\]
\end{proposition}

考虑使用傅里叶变换的伸缩性质，可以得到：
\begin{proposition}
    设$\mathcal{L} =A(\mathbb{Z}^2)$为$\mathbb{R}^2$中的格，则由该格生成的二维狄拉克梳状分布的傅里叶变换为
    \[\mathcal{F} \shah_{\mathcal{L} }=\frac{1}{|\det A|}\shah_{A^{-T}(\mathbb{Z}^2)}\]
\end{proposition}

可以看到，空间域上，线性变换$A:\mathbb{R}^2\to\mathbb{R}^2$作用于格$\mathbb{Z}^2$，对应频域上，线性变换$A^{-T}:\mathbb{R}^2\to\mathbb{R}^2$作用于格$\mathbb{Z}^2$。这启发我们定义：
\begin{definition}[对偶格]
    设$\mathcal{L} =A(\mathbb{Z}^2)$为$\mathbb{R}^2$中的格，则由线性变换$A^{-T}$作用于$\mathbb{Z}^2$生成的格称为$\mathcal{L} $的\textbf{对偶格}(dual lattice)，记为$\mathcal{L} ^*$，即
    \[\mathcal{L} ^*=A^{-T}(\mathbb{Z}^2)\]
    在这种记号下，二维狄拉克梳状分布的傅里叶变换可以记为
    \[\mathcal{F} \shah_{\mathcal{L} }=\frac{1}{\mathrm{Area}(\mathcal{L} )}\shah_{\mathcal{L} ^*}\]

    对偶格的等价定义为：
    \[\mathcal{L} ^*=\{\mathbf{y}\in\mathbb{R}^2\left|\mathbf{y}\cdot \mathbf{x}\in\mathbb{Z},\forall \mathbf{x}\in\mathcal{L}\right. \}\]
\end{definition}
\begin{definition}[对偶基与对偶线性泛函]
    考虑对偶格的基向量。令
    \[A=\begin{bmatrix}
        \mathbf{u}_1&\mathbf{u}_2
    \end{bmatrix},A^{-T}=\begin{bmatrix}
        \mathbf{u}^*_1&\mathbf{u}^*_2
    \end{bmatrix}\]
    其中$\mathbf{u}^*_1,\mathbf{u}^*_2$为对偶格$\mathcal{L} ^*$的基向量，它们满足：
    \[A^{-1}A=\begin{bmatrix}
        \mathbf{u}^*_1&\mathbf{u}^*_2
    \end{bmatrix}^T\begin{bmatrix}
        \mathbf{u}_1&\mathbf{u}_2
    \end{bmatrix}=I=\begin{bmatrix}
        1&0\\
        0&1
    \end{bmatrix}\]
    即\[\mathbf{u}^*_i\cdot \mathbf{u}_j=\begin{cases}
    1,&\mathrm{i}=j\\
    0,&\mathrm{i}\neq j
    \end{cases}=\delta^i_{j}\]
    $\delta^i_{j}$为克罗内克δ符号。我们将$\mathbf{u}^*_1,\mathbf{u}^*_2$称为$\mathbf{u}_1,\mathbf{u}_2$的\textbf{对偶基}。\textbf{对偶基生成了对偶格}。此外，与对偶基做内积，定义了内积空间上唯一一组满足正交归一化条件，即$f_i(\mathbf{u_j})=\delta^i_{j}$的线性泛函，这组线性泛函称为\textbf{对偶线性泛函}。
\end{definition}

$\shah_{\mathcal{L} }$的$\delta$分布位于格$\mathcal{L} $的格点上，而其傅里叶变换$\mathcal{F} \shah_{\mathcal{L} }$的$\delta$分布位于对偶格$\mathcal{L} ^*$的格点上。这为我们展示了一种理解空间域与频域关系的方式，并从更高的视角解释了我们在\ref{sub:fourier_of_dirac_comb}最后所发现的现象，即一维狄拉克梳状分布经过傅里叶变换，$\delta$的位置具有“倒数”关系。

最后，我们来讨论二维信号的采样与重建问题。
\begin{definition}[二维带限函数]
    设$f\in L^1(\mathbb{R}^2)$，如果存在矩形区域$$B=\{( \xi_1,\xi_2)\left| |\xi_1|\leq \xi_{m1},|\xi_2|\leq \xi_{m2}\right.\}$$
    使得$\text{supp} \mathcal{F}f \subset B$，就说$f$为二维\textbf{带限函数}(bandlimited function)，带宽为$(\xi_{m1},\xi_{m2})$。
\end{definition}
特别地，我们考虑
$$\text{supp}\mathcal{F} f\subset\left\{( \xi_1,\xi_2)\left| |\xi_1|\leq \dfrac{1}{2},|\xi_2|\leq \dfrac{1}{2}\right.\right\}$$
的情况。由恒等式
\[\mathcal{F} f(\boldsymbol{\xi})=\chi_{[-1/2,1/2]^2}(\boldsymbol{\xi})(\mathcal{F} f(\boldsymbol{\xi})*\shah)\]
取傅里叶逆变换，由\ref{sub:multi_fourier}中的结果可知$\mathcal{F} ^{-1}\chi_{[-1/2,1/2]}(\mathbf{x})=\text{sinc}(\mathbf{x})$，我们得到：\begin{align*}
    f(\mathbf{x})&=\mathcal{F} ^{-1}[\chi_{[-1/2,1/2]^2}(\boldsymbol{\xi})(\mathcal{F} f(\boldsymbol{\xi})*\shah)](\mathbf{x})\\
    &=\mathcal{F} ^{-1}\chi_{[-1/2,1/2]^2}(\mathbf{x})*(f(\mathbf{x})*\shah(\mathbf{x}))\\
    &=\text{sinc}(\mathbf{x})*(f(\mathbf{x})*\shah(\mathbf{x}))\\
    &=\sum_{\mathbf{k}\in\mathbb{Z}^2}f(\mathbf{k})\text{sinc}(x_1-k_1)\text{sinc}(x_2-k_2)
\end{align*}
最后一行用到了卷积的结合律，见\ref{sub:properties_of_convolution}。综上得到以下引理：
\begin{lemma}[方形采样公式]
    设$$f\in L^1(\mathbb{R}^2),\text{supp}\mathcal{F} f\subset\left\{( \xi_1,\xi_2)\left| |\xi_1|\leq \dfrac{1}{2},|\xi_2|\leq \dfrac{1}{2}\right.\right\}$$
    则有
    \[f(\mathbf{x})=\sum_{\mathbf{k}\in\mathbb{Z}^2}f(\mathbf{k})\text{sinc}(x_1-k_1)\text{sinc}(x_2-k_2)\]
\end{lemma}

更一般地，我们考虑$\text{supp}\mathcal{F} f$属于格$\mathcal{L} =A(\mathbb{Z}^2),A=[\mathbf{u}_1,\mathbf{u}_2]$的某个基本平行四边形内。显而易见，带限等价于存在这样一个平行四边形，使得频域信号$\mathcal{F} f$支于它。由于我们通常考虑实信号，不妨设这个基本平行四边形是关于原点对称的。

我们希望利用已经得到的方形采样公式来推导格采样公式。为此，考虑将频域信号$\mathcal{F} f$进行伸缩变换，这对应时域上“对偶”的伸缩变换：
\[\mathcal{F} f(A\boldsymbol{\xi})=\frac{1}{|\det A|}\mathcal{F}[f(A^{-T}\mathbf{x})](\boldsymbol{\xi})\]

令$g(\mathbf{x})=f(A^{-T}\mathbf{x})$，则$\text{supp}\mathcal{F} g\in[-1/2,1/2]^2$，因此可以应用方形采样公式：
\[g(\mathbf{x})=\sum_{\mathbf{k}\in\mathbb{Z}^2}g(\mathbf{k})\text{sinc}(x_1-k_1)\text{sinc}(x_2-k_2)\]
将$g$的定义代入上式，得到：
\[f(A^{-T}\mathbf{x})=\sum_{\mathbf{k}\in\mathbb{Z}^2}f(A^{-T}\mathbf{k})\text{sinc}(x_1-k_1)\text{sinc}(x_2-k_2)\]

以上结果可以改写为\textbf{格采样公式}：
\begin{theorem}[格采样公式]
    设$$f\in L^1(\mathbb{R}^2),\text{supp}\mathcal{F} f\subset\left\{( \xi_1,\xi_2)\left| \mathbf{\xi}=A\boldsymbol{\eta},\boldsymbol{\eta}\in[-1/2,1/2]^2\right.\right\}$$
    则有
    \[f(\mathbf{y})=\sum_{k_1,k_2\in\mathbb{Z}}f(k_1\mathbf{u}^*_1+k_2\mathbf{u}^*_2)\text{sinc}(\mathbf{u}_1\cdot \mathbf{y}-k_1)\text{sinc}(\mathbf{u}_2\cdot \mathbf{y}-k_2)\]
\end{theorem}
\begin{proof}
    令$\mathbf{y}=A^{-T}\mathbf{x}$，则$\mathbf{x}=A^T\mathbf{y}$，于是
    \[x_1=\mathbf{u}_1\cdot \mathbf{y},x_2=\mathbf{u}_2\cdot \mathbf{y}\]
    对$A^{-T}\mathbf{k}$，我们知道$A^{-T}$的行向量分别为$\mathbf{u}^*_1,\mathbf{u}^*_2$，因此$A^{-T}\mathbf{k}=k_1\mathbf{u}^*_1+k_2\mathbf{u}^*_2$。综上得到：
    \begin{align*}
        f(\mathbf{y})&=\sum_{\mathbf{k}\in\mathbb{Z}^2}f(A^{-T}\mathbf{k})\text{sinc}(\mathbf{u}_1\cdot \mathbf{y}-k_1)\text{sinc}(\mathbf{u}_2\cdot \mathbf{y}-k_2)\\
        &=\sum_{k_1,k_2\in\mathbb{Z}}f(k_1\mathbf{u}^*_1+k_2\mathbf{u}^*_2)\text{sinc}(\mathbf{u}_1\cdot \mathbf{y}-k_1)\text{sinc}(\mathbf{u}_2\cdot \mathbf{y}-k_2)
    \end{align*}
\end{proof}

称之为格采样公式，是因为采样点位于对偶格$\mathcal{L} ^*=A^{-T}(\mathbb{Z}^2)$上。下面左图绘制了空间域上的采样点，而右图则绘制了频域上的带限区域。频域上的限制越强，频域上的格的基本平行四边形就越小，基向量越短，这样，空间域上的格就越大，基向量越长，节点也越稀疏，即重建信号所需的采样点就越少。这符合我们的直观认识。
\begin{figure}[H]
    \centering
    \includegraphics[width=0.8\textwidth]{lattice_sample}
    \caption{格采样示意图}
\end{figure}
